% 05_4 扩展分析（中文）
\paragraph{分数重叠与标定策略。}
真实 ShezhenV3 验证对存在分数重叠：清晰组的最小 $Q_{\mathrm{img}}$ 可能低于模糊组的最大 $Q_{\mathrm{img}}$。
原因是舌象纹理与苔质在轻度失焦下仍贡献 Laplacian 能量。
因此同时报告组中位数与 Spearman $\rho$，不假设完全可分簇。
$\tau$ 取两组中位数的中点，并限制在 $[0.45,0.65]$；协议与阶段 2 一致，但以 572 张真实清晰图及成对模糊替代合成直方图。

\paragraph{部署含义。}
尽管存在重叠，B2 在测试矩阵上的 M2 仍为 0.604（物理重采，3 seeds），因为闭环重采与真实分数结合，仍可在 $K{=}2$ 内越过 $\tau$。
运维人员可查阅 JSONL 中的 \texttt{flags} 诊断失败原因，无需检查原始云上传内容。

\paragraph{与合成阶段 2 的对比。}
合成标定：$\tau{=}0.505$，中位数 $0.819/0.191$；ShezhenV3：$\tau{=}0.498$，中位数 $0.564/0.431$，M5 $\rho{=}0.47$（合成约 0.75）。
真实纹理使 $Q_{\mathrm{img}}$ 整体下移，但 B2 相对 B0/B1 的路由收益在测试集上保持，支持架构在玩具数据之外的外部有效性。
